\documentclass[11pt,a4paper]{article}
\usepackage[T1]{fontenc}
\usepackage[utf8]{inputenc}
\usepackage[margin=23mm,headheight=15pt]{geometry}
\usepackage{amsmath,amssymb,mathtools,amsthm}
\usepackage{newpxtext,newpxmath}
\usepackage{microtype,xcolor,enumitem,environ,needspace,fancyhdr,etoolbox}
\usepackage[colorlinks=true,linkcolor=heading,urlcolor=heading]{hyperref}
\definecolor{heading}{HTML}{244B65}
\definecolor{hintcolor}{HTML}{765638}
\setlength{\parindent}{0pt}
\setlength{\parskip}{5pt plus 1pt}
\setlength{\emergencystretch}{2em}
\setlist{topsep=4pt,itemsep=4pt,parsep=2pt,leftmargin=*}
\allowdisplaybreaks[1]
\newtheoremstyle{exostyle}{8pt}{6pt}{\normalfont}{}{\color{heading}\bfseries}{.}{\newline}{}
\theoremstyle{exostyle}
\newtheorem{exercise}{Exercise}
\BeforeBeginEnvironment{exercise}{\par\addvspace{10pt}\Needspace{9\baselineskip}%
  \begingroup\color{heading!45}\hrule height .6pt\endgroup\nobreak}
\newif\ifcorr
\ifdefined\StatementOnly\corrfalse\else\corrtrue\fi
\newcommand{\showsolutions}{\corrtrue}
\newcommand{\hidesolutions}{\corrfalse}
\NewEnviron{hints}{\ifcorr\par\Needspace{4\baselineskip}\noindent
  {\color{hintcolor}\bfseries Useful hints}\par\BODY\par\medskip\fi}
\NewEnviron{solution}{\ifcorr\par\Needspace{4\baselineskip}\noindent
  {\color{heading}\bfseries Detailed solution}\par\BODY\par\medskip\fi}
\newcommand{\R}{\mathbb{R}}
\newcommand{\push}{\#}
\newcommand{\W}{\mathcal{W}}
\newcommand{\Wone}{\mathcal{W}_1}
\newcommand{\Wtwo}{\mathcal{W}_2}
\newcommand{\Unif}{\mathcal U}
\newcommand{\ip}[2]{\langle #1,#2\rangle}
\newcommand{\gradW}{\mathord{\nabla\mkern-9.3mu\nabla}}
\pagestyle{fancy}
\fancyhf{}
\fancyhead[L]{\small\color{heading}Optimal Transport / Exercises}
\fancyhead[R]{\small\ifcorr Solutions\else Statements\fi}
\fancyfoot[C]{\small\thepage}
\renewcommand{\headrulewidth}{0.3pt}
\renewcommand{\maketitle}{\thispagestyle{plain}\begingroup
  \color{heading}\LARGE\bfseries\SheetTitle\par\endgroup
  \vspace{3pt}{\large\ifcorr Statements, hints and worked solutions\else Exercise statements\fi}\par
  \vspace{7pt}\hrule\vspace{9pt}
  \ifcorr Each exercise is followed by question-by-question hints and a worked solution.
  Try the questions before reading the hints.\else
  This version contains the exercises only. Hints and worked solutions are available in the companion PDF.\fi
  \par\medskip}
\hypersetup{pdfauthor={Optimal Transport for Machine Learners},pdfsubject={Exercises in optimal transport}}

\newcommand{\SheetTitle}{Wasserstein Distances and Duality}
\newcommand{\X}{\mathcal X}
\newcommand{\Y}{\mathcal Y}
\newcommand{\U}{\mathcal U}
\hypersetup{pdftitle={Wasserstein Distances and Duality -- Exercises},
  pdfauthor={Julie Delon (original exercises); OT4ML (adaptation and solutions)}}

\begin{document}
\maketitle
\textbf{Original exercises: Julie Delon.} Adapted for OT4ML.

\textbf{Notation.}
Write $\U(\alpha,\beta)$ for probability couplings with marginals
$\alpha,\beta$, and $\mathcal L_c(\alpha,\beta)$ for the infimum of
$\int c\,d\pi$ over these couplings. We use the full total variation
norm $\|\nu\|_{\mathrm{TV}}=|\nu|(\X)$ for signed measures; for example,
$\|\sum_i v_i\delta_{x_i}\|_{\mathrm{TV}}=\sum_i|v_i|$ for distinct points.

\begin{exercise}[Total variation and the $0/1$ cost]\label{ex-dd-tv}
On a standard Borel space $\X$, set $d_0(x,y)=\mathbf 1_{\{x\ne y\}}$.
\begin{enumerate}[label=(\alph*)]
\item Let $\alpha=\sum_i a_i\delta_{x_i}$ and
  $\beta=\sum_i b_i\delta_{x_i}$ be probabilities on the same finite
  set of distinct points (zero weights are allowed). Compute
  $\mathcal L_{d_0}(\alpha,\beta)$ and construct an optimal plan.
\item Extend the identity $\mathcal L_{d_0}(\alpha,\beta)=
  \tfrac12\|\alpha-\beta\|_{\mathrm{TV}}$ to arbitrary probabilities on
  $\X$. Explain why, on a countable discrete space, this also gives
  $\mathcal W_{p,d_0}(\alpha,\beta)^p=\tfrac12\|\alpha-\beta\|_{\mathrm{TV}}$
  for $1\le p<\infty$, where the subscript $d_0$ specifies the ground metric.
\end{enumerate}
\end{exercise}
\begin{hints}
\begin{enumerate}[label=(\alph*)]
\item The cost of a plan is the mass it sends off the diagonal.
  How much mass can you keep at each $x_i$? Couple the two residual
  weight vectors after taking out that mass.
\item Use $\lambda=\alpha+\beta$ as a dominating measure and take the
  pointwise minimum of the two densities. For any plan, the measure
  of its diagonal part is dominated by both marginals.
\end{enumerate}
\end{hints}
\begin{solution}
\textbf{(a) Match as much mass as possible in place.}
A coupling matrix $\mathrm P$ has nonnegative entries, row sums $a_i$
and column sums $b_j$. Its cost is
\[
  \sum_{i,j}\mathbf 1_{\{i\ne j\}}\mathrm P_{i,j}
    =1-\sum_i\mathrm P_{i,i}.
\]
The row and column constraints imply separately that
$\mathrm P_{i,i}\le a_i$ and $\mathrm P_{i,i}\le b_i$.
Set $m_i=\min(a_i,b_i)$ and $r=1-\sum_i m_i$. No plan can cost less than
$r$, and the identity $|a_i-b_i|=a_i+b_i-2m_i$ gives
\[
  r=\frac12\sum_i|a_i-b_i|
    =\frac12\|\alpha-\beta\|_{\mathrm{TV}}.
\]
To attain the lower bound, keep $m_i$ at each $x_i$. The residual weights
$a_i'=a_i-m_i$ and $b_i'=b_i-m_i$ both sum to $r$ and satisfy
$a_i'b_i'=0$ for every $i$. If $r>0$, set
\[
  \mathrm P^\star_{i,j}=\mathbf 1_{\{i=j\}}m_i+\frac{a_i'b_j'}{r}.
\]
Indeed, $\sum_j\mathrm P^\star_{i,j}=m_i+a_i'=a_i$, and
$\sum_i\mathrm P^\star_{i,j}=m_j+b_j'=b_j$. The product term
has no diagonal entries and has total mass
$(\sum_i a_i')(\sum_j b_j')/r=r$, so the cost is exactly $r$.
If $r=0$, the two weight vectors are equal and the diagonal plan is
feasible with cost zero. In either case the lower bound is attained;
when $r>0$, there can be other ways to match the leftover mass.

\par\Needspace{6\baselineskip}
\textbf{(b) The common part of two measures.}
Let $a=d\alpha/d\lambda$ and $b=d\beta/d\lambda$ for
$\lambda=\alpha+\beta$. Define the common part $\eta$ by
$d\eta=\min(a,b)\,d\lambda$. The residuals $\alpha'=\alpha-\eta$
and $\beta'=\beta-\eta$ have densities $(a-b)_+$ and $(b-a)_+$.
They are concentrated on the disjoint measurable sets $\{a>b\}$
and $\{b>a\}$, respectively, and their masses are both
\[
  r=1-\eta(\X)=\frac12\int |a-b|\,d\lambda
               =\frac12\|\alpha-\beta\|_{\mathrm{TV}}.
\]
Equivalently, $r=\sup_{A}|\alpha(A)-\beta(A)|$, where the supremum is
over measurable sets: the positive and negative parts of $\alpha-\beta$
each have mass $r$, and $A=\{a>b\}$ attains it. This is why the
\emph{probability distance} often called total variation is half the
signed-measure norm used here.

Let $\Delta(x)=(x,x)$ and let $D=\{(x,x):x\in\X\}$.
For a coupling $\pi$, its diagonal mass defines a measure
$\gamma(A)=\pi(D\cap(A\times\X))$ on $\X$. The diagonal is measurable
on a standard Borel space (a countable separating family of measurable
sets detects whether two points agree). Since
$D\cap(A\times\X)=D\cap(\X\times A)$, the marginal constraints give
$\gamma(A)\le\alpha(A)$ and $\gamma(A)\le\beta(A)$.
Hence $\gamma\ll\lambda$ and its density is at most both $a$ and $b$,
so $\gamma\le\eta$. It follows that
\[
  \int d_0\,d\pi=1-\pi(D)=1-\gamma(\X)\ge 1-\eta(\X)=r .
\]
For $r>0$, take $\pi^\star=\Delta_\push\eta+\alpha'\otimes\beta'/r$.
The first and second marginals of the product term are $\alpha'$ and
$\beta'$ because the opposite residual has mass $r$. Adding the
diagonal term therefore restores the two prescribed marginals.
The product term is off the diagonal by disjointness of the residual
sets, so $\pi^\star$ costs exactly $r$. If $r=0$, use $\Delta_\push\alpha$.

Since $d_0^p=d_0$ for every finite $p\ge 1$, the same cost appears in
the definition of $\mathcal W_{p,d_0}^p$ on a countable discrete space.
On an uncountable standard Borel space the calculation still holds
for the measurable cost on the given Borel structure; it does not
require changing that structure to the discrete metric topology.
\end{solution}

\begin{exercise}[Equivalence of topologies on a finite metric space]\label{ex-dd-finite}
Let $(\X,d)$ be a finite metric space with at least two points, and set
$d_{\min}=\min_{x\ne y}d(x,y)$ and $d_{\max}=\max_{x,y}d(x,y)$.
\begin{enumerate}[label=(\alph*)]
\item Prove that all probability measures $\alpha,\beta$ on $\X$ satisfy
  \[
    \frac{d_{\min}}2\|\alpha-\beta\|_{\mathrm{TV}}
    \le \Wone(\alpha,\beta)
    \le\frac{d_{\max}}2\|\alpha-\beta\|_{\mathrm{TV}} .
  \]
\item Deduce that $\Wone$ and total variation induce the same topology.
\end{enumerate}
\end{exercise}
\begin{hints}
\begin{enumerate}[label=(\alph*)]
\item Bound $d$ above and below by multiples of $d_0$. For the upper
  bound, test with the optimal $d_0$ plan from Exercise~\ref{ex-dd-tv}.
\item Turn the two distance estimates into ball inclusions at a fixed
  center. Why is $d_{\min}>0$?
\end{enumerate}
\end{hints}
\ifcorr\Needspace{12\baselineskip}\fi
\begin{solution}
\textbf{(a) Price the mass that must move.}
On the diagonal both $d$ and $d_0$ vanish. Off the diagonal $d_0=1$,
so the definitions of $d_{\min}$ and $d_{\max}$ give
$d_{\min}d_0(x,y)\le d(x,y)\le d_{\max}d_0(x,y)$ everywhere.
For any coupling $\pi$, the lower bound and Exercise~\ref{ex-dd-tv} give
\[
  \int d\,d\pi\ge d_{\min}\int d_0\,d\pi
                \ge \frac{d_{\min}}2\|\alpha-\beta\|_{\mathrm{TV}}.
\]
Taking the infimum proves the first inequality. For the second use the
specific coupling $\pi^\star$ that is optimal for $d_0$. It need not be
optimal for $d$:
\[
  \Wone(\alpha,\beta)\le\int d\,d\pi^\star
    \le d_{\max}\int d_0\,d\pi^\star
    =\frac{d_{\max}}2\|\alpha-\beta\|_{\mathrm{TV}} .
\]
The constants are sharp. For $\alpha=\delta_x$, $\beta=\delta_y$ with
$x\ne y$, the TV norm of the difference is $2$ and the unique coupling
has cost $d(x,y)$. Choosing pairs attaining $d_{\min}$ and $d_{\max}$
gives equality in the respective bounds.

\textbf{(b) Compare neighborhoods.}
Finiteness and definiteness of $d$ give $0<d_{\min}\le d_{\max}<\infty$.
Writing $B_{\mathrm{TV}}$ for balls in the full TV norm and $B_{\Wone}$
for Wasserstein balls, part~(a) gives, for every center $\alpha$
and radius $\varepsilon>0$,
\[
  B_{\mathrm{TV}}(\alpha,2\varepsilon/d_{\max})
   \subset B_{\Wone}(\alpha,\varepsilon)
   \subset B_{\mathrm{TV}}(\alpha,2\varepsilon/d_{\min}).
\]
Each kind of neighborhood therefore contains one of the other kind.
In particular, $\Wone(\alpha_n,\alpha)\to0$ if and only if
$\|\alpha_n-\alpha\|_{\mathrm{TV}}\to0$. If $\X$ has only one point,
the two topologies coincide trivially and no $d_{\min}$ is needed.
\end{solution}

\begin{exercise}[Algebra of $c$-transforms]\label{ex-dd-transforms}
Let $\X,\Y$ be nonempty compact metric spaces and
$c:\X\times\Y\to\R$ continuous. For continuous real-valued potentials
$f,f'$ on $\X$ and $g$ on $\Y$, use the conventions
\[
  f^c(y)=\inf_{x\in\X}\{c(x,y)-f(x)\},\qquad
  g^{\bar c}(x)=\inf_{y\in\Y}\{c(x,y)-g(y)\}.
\]
Here $\bar c(y,x)=c(x,y)$, and the transforms are applied in the order
written: $f^{c\bar c}=(f^c)^{\bar c}$.
Prove the following pointwise statements:
\begin{enumerate}[label=(\alph*)]
\item If $f\le f'$, then $f^c\ge(f')^c$.
\item $f^{c\bar c}\ge f$.
\item $g^{\bar c c}\ge g$.
\item $f^{c\bar c c}=f^c$.
\end{enumerate}
\end{exercise}
\begin{hints}
\begin{enumerate}[label=(\alph*)]
\item Compare the expressions inside the infima before taking them.
\item In $f^c(y)$, test the infimum with a fixed $x$; then rearrange.
\item Repeat the argument with the roles of $\X$ and $\Y$ exchanged.
\item For one inequality, apply (a) to (b). For the other, plug
  $g=f^c$ into (c).
\end{enumerate}
\end{hints}
\ifcorr\Needspace{12\baselineskip}\fi
\begin{solution}
Compactness and continuity imply that the transforms are finite.
For example,
$|f^c(y)-f^c(y')|\le\sup_x |c(x,y)-c(x,y')|$: the two functions
of $x$ inside the infima differ by at most the right-hand side.
Uniform continuity of $c$ on the compact product gives continuity of
$f^c$. The same reasoning applies to $g^{\bar c}$, so the repeated
transforms are also well defined.

\textbf{(a) A single transform reverses the order.}
If $f(x)\le f'(x)$ for all $x$, then
$c(x,y)-f(x)\ge c(x,y)-f'(x)$ for each $x,y$.
Taking the infimum over $x$ on both sides preserves this last
inequality, giving $f^c\ge(f')^c$. The same proof gives order reversal
for the $\bar c$-transform.

\textbf{(b) Closing a potential on the source.}
By definition $f^c(y)\le c(x,y)-f(x)$ for every $x,y$. Thus for each
fixed $x$, all the values $c(x,y)-f^c(y)$ are at least $f(x)$.
Taking their infimum gives
\[
  f^{c\bar c}(x)=\inf_{y\in\Y}\{c(x,y)-f^c(y)\}\ge f(x).
\]
The starting inequality also says that $(f,f^c)$ is dual feasible:
$f(x)+f^c(y)\le c(x,y)$ everywhere.

\textbf{(c) Closing a potential on the target.}
This time $g^{\bar c}(x)\le c(x,y)-g(y)$ for every $x,y$.
Rearrange to $c(x,y)-g^{\bar c}(x)\ge g(y)$, and take the infimum
over $x$ for each fixed $y$. This gives $g^{\bar c c}(y)\ge g(y)$.

\textbf{(d) The third transform repeats the first.}
From (b), $f\le f^{c\bar c}$, so (a) gives
$f^c\ge(f^{c\bar c})^c=f^{c\bar c c}$.
On the other hand, apply (c) to the target potential $g=f^c$ to
obtain $(f^c)^{\bar c c}\ge f^c$. The left-hand side is
$f^{c\bar c c}$, proving equality. In particular, applying $\bar c$
to that equality gives $f^{c\bar c c\bar c}=f^{c\bar c}$: taking the
double transform twice gives the same result as taking it once.

\par\Needspace{5\baselineskip}
\textbf{Consequence for alternating dual updates.}
With the other potential fixed, a transform is the pointwise
largest feasible choice. Starting from $f$, updating first the target
potential and then the source gives the feasible pair
$(f^{c\bar c},f^c)$. By (d), the next target update returns $f^c$ again,
and the next source update returns $f^{c\bar c}$. This proves that
the updates stop, not that the pair is jointly optimal.

For instance, on two points with diagonal costs $0$ and off-diagonal
costs $1$, start at $f=(0,1)$. The two columnwise minima give
$f^c=(0,-1)$, and the two rowwise minima give $f^{c\bar c}=(0,1)$.
For source weights $(3/4,1/4)$ and target weights $(1/4,3/4)$, this
fixed pair has dual value $1/4-3/4=-1/2$. In fact the optimal value
is $1/2$: keep mass $1/4$ at each point and move the remaining $1/2$
from the first to the second. The feasible potentials $(1,0)$ and
$(-1,0)$ attain this same value. What the fixed pair misses is
equality on the off-diagonal edge that must carry mass.
\end{solution}

\par\addvspace{7pt}\Needspace{3\baselineskip}
\textbf{Further reading.}
See the discussions of total variation in \emph{Wasserstein Space}
and of the algebra of $c$-transforms in \emph{Dual Problem} in
\href{https://www.gpeyre.com/ot4ml/OT4ML/OT4ML.pdf}
     {\emph{Optimal Transport for Machine Learners}}.
\end{document}
